\section{Empirical Fuel-Corrected Tire Degradation Physics}
\label{sec:tire_physics_degradation}

Tire degradation represents one of the primary physical determinants of race pace and pit strategy viability \cite{bekker2009stochastic}. In Free Practice (FP2) long runs, teams assess compound longevity across continuous stints of 10 to 25 laps. However, naively regressing raw telemetry lap times against lap number introduces severe econometric bias due to the confounding effect of vehicle mass burnoff.

\subsection{Confounding Fuel Mass Dynamics}
During a full Grand Prix distance ($D \approx 305$~km), a modern Formula 1 vehicle consumes approximately $100$~kg to $110$~kg of fuel across $L_{\text{race}} \approx 50$--$70$ laps, yielding an average fuel consumption rate of:
\begin{equation}
\dot{m}_{\text{fuel}} \approx 1.8 \text{ to } 2.0 \text{ kg / lap}
\end{equation}
Vehicle mass reduction reduces the normal load and inertial resistance through corners and acceleration phases, producing an empirical lap time improvement:
\begin{equation}
\lambda_{\text{fuel}} \approx 0.035 \text{ seconds / lap / kg-burnoff}
\end{equation}
Consequently, if true physical tire degradation $\beta_{\text{deg}}$ is less than the fuel burnoff advantage ($\beta_{\text{deg}} < \lambda_{\text{fuel}}$), the raw observed lap times $t_{\text{raw}}(l)$ will decrease monotonically over the course of the stint, producing a spurious negative degradation slope ($\beta_{\text{raw}} < 0$). This violates physical thermodynamics and misleads strategy models into anticipating infinite tire durability.

\subsection{Fuel-Corrected Formulation}
To isolate true tire wear, the engine applies a linear fuel mass correction offset. For a stint comprising laps $l \in \{1, 2, \dots, L_{\text{stint}}\}$:
\begin{equation}
t_{\text{corrected}}(l) = t_{\text{raw}}(l) - \lambda_{\text{fuel}} \cdot (L_{\text{stint}} - l)
\label{eq:fuel_corrected_lap}
\end{equation}
where $\lambda_{\text{fuel}} = 0.035$~s/lap. Under Eq.~\eqref{eq:fuel_corrected_lap}, lap times are normalized to constant full-fuel equivalent mass.

\subsection{Telemetry Filtering and Slope Estimation}
Prior to slope estimation, raw telemetry extracted via the FastF1 interface \cite{roembke2021fastf1} is filtered according to strict green-flag telemetry integrity criteria:
\begin{enumerate}
    \item \textbf{In/Out-Lap Elimination:} Stint entry (pit exit) and in-laps (pit entry) are discarded.
    \item \textbf{Traffic and Yellow Flag Filtering:} Laps with track status flags $\ne 1$ (Safety Car, VSC, Yellow) or exceeding the session personal best by more than 7\% ($t_{\text{raw}} > 1.07 \cdot t_{\text{best}}$) are pruned.
    \item \textbf{Minimum Stint Length:} Only stints with $N_{\text{clean}} \ge 5$ consecutive valid laps are accepted.
\end{enumerate}

The physical tire degradation gradient $\beta_{\text{deg}}$ is estimated via Ordinary Least Squares (OLS):
\begin{equation}
\beta_{\text{deg}} = \frac{\sum_{l=1}^{L} (l - \bar{l})(t_{\text{corrected}}(l) - \bar{t}_{\text{corrected}})}{\sum_{l=1}^{L} (l - \bar{l})^2}
\label{eq:deg_slope}
\end{equation}
Following Eq.~\eqref{eq:deg_slope}, empirical tests across Bahrain and Barcelona telemetry confirm that $\beta_{\text{deg}} > 0$ across all dry compounds (typically $0.020$--$0.085$~s/lap for C3 Medium and $0.050$--$0.140$~s/lap for C4/C5 Soft compounds).

\subsection{Offline Scikit-Learn Regression Fallback}
When external telemetry APIs are unreachable or when operating in air-gapped evaluation environments, the engine activates the \texttt{SklearnTireDegradationModel}. This fallback module fits a Multi-Layer Perceptron (MLP) regressor to historical Pirelli compound wear curves conditioned on:
\begin{equation}
\hat{\beta}_{\text{deg}} = f_{\boldsymbol{\theta}}\left( \text{compound}, T_{\text{track}}, L_{\text{stint}}, \text{abrasion\_index} \right)
\label{eq:tire_fallback}
\end{equation}
ensuring deterministic degradation gradients without external network dependencies.
